\section{Assumed background}
\label{sec:background}

\emph{Source: handout preamble (``Assumed background: DAGs, the do-operator,
backdoor adjustment, potential outcomes, regression, a little machine
learning.'').} Everything from \S\ref{sec:bigpicture} onward leans on all six
of these.

\textbf{Running example}, used throughout: age $W \to$ exercise $A \to$
weight $M \to$ blood pressure $R$; a model $f$ predicts 10-year
cardiovascular risk, $\hat Y = f(X)$.

\subsection{Notation}
\label{sec:notation}

A quick-reference table for every symbol used in these notes --- consult it
whenever a formula reads ``obviously'' but isn't.

\begin{longtable}{@{}p{3.0cm}p{10.2cm}@{}}
\toprule
\textbf{Symbol} & \textbf{Meaning} \\
\midrule
\endhead
$W, A, M, R$ & The running example's features: age, exercise, weight,
  blood pressure (in that causal order). \\
$X$ & The full feature vector, e.g.\ $(W,A,M,R)$; $X_{-A}$ means ``all
  features except $A$.'' \\
$f$ & The fitted prediction model being explained (any learner). \\
$\hat Y$ & The model's output, added to the causal graph as a node,
  $\hat Y = f(X)$. \\
$Y^*$ & The ``constructed outcome'' $f(X)$ treated as a variable for
  causal-inference purposes --- the same quantity as $\hat Y$, written
  differently depending on whether the point being made is graphical
  ($\hat Y$) or statistical ($Y^*$). \\
$P(\cdot)$, $P(dw)$ & A probability distribution; $P(dw)$ is shorthand for
  integrating/summing against the distribution of $W$ (treat it as
  ``weighted by how $W$ is actually distributed in the population''). \\
$E[\cdot]$, $E_W[\cdot]$ & Expectation; a subscript names which variable's
  distribution the average is taken over. \\
$\operatorname{do}(A=a)$ & Pearl's do-operator: \emph{set} $A$ to $a$ by
  intervention (delete the arrow into $A$), as opposed to conditioning on
  $A=a$ in observational data (see below). \\
$\perp$ & Statistical independence, e.g.\ $\{R(a)\}\perp A \mid W$ reads
  ``independent of $A$, once you condition on $W$.'' \\
$\theta(a)$ & The population target of a dependence plot at level $a$
  (e.g.\ the average TDP, Eq.~\ref{eq:tdp-backdoor}). \\
$\hat\theta(a)$ & An \emph{estimator} of $\theta(a)$ from data. A hat over
  any symbol ($\hat\mu$, $\hat\pi$, $\hat{\mathrm{se}}$, \dots) always means
  ``estimated from a finite sample,'' as opposed to the population quantity
  underneath it. \\
$\mu(a,z)$ & The outcome regression, $E[Y^* \mid A=a, Z=z]$
  (\S\ref{sec:estimation}). \\
$\pi(a\mid z)$ & The propensity/treatment regression, $P(A=a\mid Z=z)$
  (\S\ref{sec:estimation}). \\
$Z$ & A generic adjustment/covariate set (e.g.\ the parents $W$, or a
  larger valid set) --- kept as a separate letter from the running
  example's $W,M,R$ so formulas stay general across choices of adjustment
  set. \\
$Q(dr\mid a,w)$ & The \textbf{reduced form}: the conditional law of the
  non-parent features given $(A,W)$ (\S\ref{sec:identification}). \\
$\Psi(a)$ & The \textbf{bridge mean} $E[Y(A,Z(a))]$ that the NIDP turns out
  to equal (\S\ref{sec:identification}). \\
$\sqrt{n}$ (rate) & Shorthand for ``a regular estimator whose error shrinks
  like $1/\sqrt n$'' --- the standard, non-exotic rate that licenses
  ordinary $\hat\theta \pm 1.96\,\hat{\mathrm{se}}$ confidence intervals. \\
$o_p(n^{-1/2})$ & ``Shrinks faster than $1/\sqrt n$, in probability'' ---
  the condition cross-fitting needs the \emph{product} of two nuisance
  errors to satisfy (\S\ref{sec:estimation}). \\
$g_k(a)$, $B_k$, $\psi_k$ & A bin's fixed weight function, the bin itself,
  and the resulting bin target $\psi_k = \int_{B_k}\theta(a)g_k(a)\,da$
  (\S\ref{sec:continuous}). \\
$K_h$, $\theta_h$ & A smoothing kernel at bandwidth $h$, and the resulting
  smoothed curve $\theta_h = K_h * \theta$ (\S\ref{sec:continuous}). \\
$J(a)$ & The interaction remainder in
  $\theta(a)-B = [\mathrm{NDDP}(a)-B]+[\mathrm{NIDP}(a)-B]+J(a)$
  (\S\ref{sec:bigpicture}). \\
Partial-$R^2$ & The fraction of otherwise-unexplained variance an
  unmeasured confounder $U$ would need to explain to matter --- the
  currency of the sensitivity analysis in \S\ref{sec:propagation}. \\
\bottomrule
\caption{Notation used throughout these notes.}
\label{tab:notation}
\end{longtable}

\begin{quote}
\textbf{Watch for this overload:} the handout reuses the letter $R$ for two
different things. In the running example, $R$ is \emph{blood pressure}. But
in \S\ref{sec:identification}'s general identification result, $R$ denotes
``the features other than $A$ and its parents $W$'' --- which in the
running example is actually the \emph{pair} $(M, \text{blood pressure})$
together, not blood pressure alone (weight $M$ is also a non-parent of
$A$). Read $R$ as ``blood pressure'' in running-example prose, and as ``the
generic non-parent remainder'' inside reduced-form formulas --- the handout
doesn't distinguish them notationally, so context has to.
\end{quote}

\subsection{DAGs (directed acyclic graphs)}

Nodes are variables, arrows are \emph{direct} causal effects (relative to
the other variables drawn), and ``acyclic'' means no sequence of arrows
loops back to its start. The \textbf{ECM} (explanatory causal model) in this
project is exactly such a graph over the model's features --- an assumption
brought to the explanation, not something derived from $f$ itself.

Vocabulary this project actually uses:
\begin{itemize}
  \item \textbf{Parents} of $A$: nodes with an arrow directly into $A$
    (here, just $W$).
  \item \textbf{Descendants} of $A$: everything reachable by following
    arrows \emph{out} of $A$ (here, $M$, $R$, $\hat Y$).
  \item \textbf{Non-descendants} of $A$: everything else (here, $W$, plus
    anything upstream of or unrelated to $A$) --- the exact set
    \S\ref{sec:estimation}'s covariate-selection result is about.
  \item \textbf{Backdoor path}: a path between $A$ and the outcome that
    starts with an arrow \emph{into} $A$ (i.e.\ runs through a common
    cause) --- the confounding path adjustment needs to block.
\end{itemize}

One caution worth internalizing early: conditioning on the wrong node can
\emph{create} bias rather than remove it (classically, conditioning on a
\textbf{collider} --- a node two arrowheads point into). This is why ``valid
adjustment set'' is a precise graph-theoretic notion (Pearl's backdoor
criterion, next), not just ``control for everything measured.''

\subsection{The do-operator}

Notation $\operatorname{do}(A=a)$: \emph{set} $A$ to $a$ by intervention, as
opposed to \emph{observing} $A=a$ in data that happened to be collected.
The distinction matters whenever $A$ shares a common cause with the outcome:

\begin{itemize}
  \item $P(R \mid A=a)$ --- among people who happen to exercise at level
    $a$ (contaminated by whatever else predicts exercising at that level,
    e.g.\ age).
  \item $P(R \mid \operatorname{do}(A=a))$ --- if everyone were \emph{forced}
    to exercise at level $a$ (severs whatever normally determines $A$, so no
    contamination).
\end{itemize}

Graphically, $\operatorname{do}(A=a)$ means: delete every arrow \emph{into}
$A$ (its value no longer depends on its natural causes), then let everything
else propagate forward as usual. This is exactly what the TDP diagram in
Figure~\ref{fig:four-plots} (\S\ref{sec:bigpicture}) draws as a severed edge
--- TDP and PCDP are literally this operation applied to the ECM.

\subsection{Backdoor adjustment}

\paragraph{What confounding actually is.} Age has \emph{two separate}
effects here: it affects blood pressure directly (general aging, nothing
to do with exercise), and it also affects how much people exercise (for
whatever reason --- older people might exercise more or less than younger
people). Because age drives \emph{both} exercise and blood pressure, a
naive comparison of blood pressure between high- and low-exercisers picks
up two things tangled together: (1) whatever blood pressure difference
exercise itself actually causes, and (2) a spurious difference that's
really just age showing up twice --- once through its direct effect on
blood pressure, once through its effect on who ends up exercising how
much. That second, spurious piece is \textbf{confounding}: age
\emph{confounds} the naive exercise--blood-pressure comparison.
Graphically, it's exactly a \textbf{backdoor path} (\S1.2 above): the path
exercise $\leftarrow$ age $\to$ blood pressure connects exercise and blood
pressure without exercise having to cause anything, simply because they
share a cause.

\paragraph{Why we bother with the marginal-averaging trick.} We don't want
``the correlation between exercise and blood pressure among whoever
happens to be at each exercise level'' (that's contaminated by age, per
above) --- we want ``what blood pressure would look like if the \emph{same}
population, with the \emph{same} age makeup, were compared at different
exercise levels.'' Forcing every exercise level to be evaluated against the
identical, population-wide age distribution (the $P(W=w)$ in
Eq.~\ref{eq:backdoor} below) is precisely how you strip out age's
contribution and isolate exercise's own effect. This is \emph{why} we do
the adjustment at all, not just how.

The bridge from an interventional question to something estimable from
observational data. \textbf{Backdoor criterion}: a covariate set $W$ is
\emph{valid for adjustment} if (a) it contains no descendant of $A$, and
(b) it blocks every backdoor path from $A$ to the outcome.

\paragraph{What ``valid'' buys you, and what goes wrong without it.}
``Valid'' means exactly this: plug $W$ into the formula below, and you get
the \emph{correct} causal quantity $P(R\mid\operatorname{do}(A=a))$ ---
not just some number, the right one. Each half of the criterion guards
against a specific way of getting a \emph{wrong} number instead:
\begin{itemize}
  \item \textbf{Violating (b)} --- a backdoor path left unblocked --- means
    some confounding is still in the room. If there were a second
    confounder, say socioeconomic status $S$, that also affects both
    exercise and blood pressure, and you adjusted only for age (not $S$),
    the path exercise $\leftarrow S \to$ blood pressure would still be
    open. The ``adjusted'' number would still be contaminated by $S$,
    exactly as in the naive calculation above, just by a different amount.
  \item \textbf{Violating (a)} --- including a descendant of $A$ --- can go
    wrong two ways. If the descendant is something $A$ causes on the way
    to the outcome (a \emph{mediator}, like weight $M$, when what you want
    is the \emph{total} effect), conditioning on it blocks part of the very
    effect you're trying to measure --- you'd end up computing something
    closer to a direct effect than a total one. If the descendant is a
    \textbf{collider} (something caused by both $A$ and something else),
    conditioning on it can manufacture a spurious association that wasn't
    there to begin with.
\end{itemize}
When both (a) and (b) hold, the formula
\begin{equation}
  P(R \mid \operatorname{do}(A=a)) \;=\; \sum_w P(R \mid A=a, W=w)\, P(W=w)
  \label{eq:backdoor}
\end{equation}
is \emph{guaranteed} to equal the true interventional quantity --- that
guarantee is the entire content of ``valid.''

\emph{In words:} to get blood pressure's distribution under a forced
exercise level $a$, go age group by age group ($w$), look at blood pressure
among the people who actually exercised at level $a$ within that age group
($P(R\mid A=a,W=w)$), and then combine those age-group numbers into one
overall number by weighting each age group by how common it actually is in
the population ($P(W=w)$) --- not by how common it is among people who
exercise at level $a$.

The crucial move: average over the \textbf{marginal} distribution of $W$
(the actual population mix), not the distribution of $W$ \emph{among people
at $A=a$} --- that second thing is what a naive $E[R\mid A=a]$ implicitly
does, and it is exactly where confounding sneaks back in.

\paragraph{Worked toy version of the running example.} Say age is 50\% young
/ 50\% old \emph{in the population overall} --- i.e.\ combining the
high-exercise and low-exercise groups together, not within each group
separately. (In real data the age mix \emph{inside} the high-exercise group
could easily be, say, 70\% old / 30\% young if older people tend to
exercise more --- that mismatch between each group's own age mix and the
population's age mix is precisely the confounding the adjustment below is
designed to remove.) And (invented purely to illustrate) blood pressure is
as in Table~\ref{tab:toy-backdoor}.

\begin{table}[h]
\centering
\begin{tabular}{lcc}
\toprule
 & high exercise & low exercise \\
\midrule
young & 110 & 120 \\
old   & 130 & 150 \\
\bottomrule
\end{tabular}
\caption{Toy blood-pressure cell means by age and exercise level.}
\label{tab:toy-backdoor}
\end{table}

Adjusted:
\begin{align*}
  \text{high exercise} &: 0.5(110) + 0.5(130) = 120 \\
  \text{low exercise}  &: 0.5(120) + 0.5(150) = 135
\end{align*}
\emph{In words:} for ``high exercise,'' take the blood pressure of a young
high-exerciser (110) and the blood pressure of an old high-exerciser (130),
and average them \emph{using the population's actual 50/50 age split} ---
not whatever age split happens to show up among the high-exercisers in the
data. Do the same for ``low exercise.'' This gives a clean $-15$ causal contrast, because both numbers are averaged over
the \emph{same} 50/50 age split. A naive $E[BP \mid A{=}\text{high}]$ vs.\
$E[BP \mid A{=}\text{low}]$ would instead average each group over
\emph{its own} age mix (however age happens to correlate with exercise in
the data).

\paragraph{Making the naive number concrete.} Suppose that, because older
people happen to exercise less in this data, the high-exercise group is
actually 80\% young / 20\% old, and the low-exercise group is 20\% young /
80\% old (this is the confounding: each group's own age mix differs from
the population's 50/50). The naive, unadjusted comparison uses \emph{each
group's own} age mix:
\begin{align*}
  \text{naive } E[BP\mid A{=}\text{high}] &= 0.8(110) + 0.2(130) = 114 \\
  \text{naive } E[BP\mid A{=}\text{low}]  &= 0.2(120) + 0.8(150) = 144
\end{align*}
giving a naive contrast of $114-144=-30$ --- \emph{twice as large} as the
true $-15$ causal effect computed above. The extra $-15$ is pure
confounding: high-exercisers look even better than they really are (via
exercise) partly because they also skew younger, and younger people have
lower blood pressure regardless of how much they exercise. Adjustment
(forcing both groups to be evaluated at the same 50/50 age split) removes
exactly that extra piece, leaving only exercise's own effect.

\paragraph{Why parents specifically?} The set of all parents of $A$ always
satisfies the backdoor criterion: any path into $A$ must pass through one of
$A$'s parents immediately before reaching it, so conditioning on all parents
blocks every such path. That is why the handout's core formula
($\theta(a) = E_W[E[Y^* \mid A=a, W]]$, $W = $ parents of $A$; see
\S\ref{sec:bigpicture}) can use \emph{just} the parents --- and why
\S\ref{sec:estimation} can then ask whether a \emph{larger} valid set (all
non-descendants) does better on variance, without needing a different
formula.

One distinction worth keeping crisp: \textbf{valid is not the same as
best.} Multiple different sets can each be valid (each plugs into
Eq.~\ref{eq:backdoor} and gives the exact correct number) while differing
in how \emph{precisely} that number can be estimated from a finite sample
--- that's the ``efficient set'' question in \S\ref{sec:estimation}, a
separate question from validity.

\subsection{Potential outcomes}

An equivalent (Neyman--Rubin) way of stating the same thing in terms of
per-person counterfactuals rather than graph surgery. For person $i$, posit
$R_i(a)$ = ``the blood pressure person $i$ would have had if their exercise
had been set to $a$'' --- only the value at their \emph{actual} $A_i$ is ever
observed; the rest are counterfactual.

Assumptions this needs (same substance as the backdoor criterion, different
language):
\begin{itemize}
  \item \textbf{Consistency}: $R_i(A_i) = R_i$ --- the potential outcome at
    the observed treatment equals what was actually seen.
  \item \textbf{Ignorability / unconfoundedness}: $\{R(a)\} \perp A \mid W$
    --- treatment is ``as good as randomly assigned'' once conditioned on
    $W$.
  \item \textbf{Positivity / overlap}: every stratum of $W$ has some chance
    of any $A$-level occurring in the data, $0 < P(A=a\mid W=w) < 1$ --- what
    \S\ref{sec:propagation}'s ``non-overlap ratio'' diagnostic checks.
\end{itemize}

Why this project needs it beyond the simple $R(a)$: the \textbf{NIDP} is a
``bridge mean'' $\Psi(a) = E[Y(A, Z(a))]$ --- a \emph{nested} potential
outcome, mixing a person's actual $A$ with the downstream features $Z$ set to
where they would be under a hypothetical $A=a$. That is genuinely
potential-outcomes language, not something the do-operator alone writes down
as easily --- see \S\ref{sec:identification}.

\subsection{Regression}

The tool that actually gets fit to data once identification says \emph{what}
to estimate. Two regressions do essentially all the work in
\S\ref{sec:estimation}:
\begin{itemize}
  \item \textbf{Outcome regression} $\mu(a,z) = E[Y^* \mid A=a, Z=z]$.
  \item \textbf{Propensity / treatment regression} $\pi(a\mid z) = P(A=a\mid Z=z)$.
\end{itemize}
AIPW combines both (and is right if \emph{either} is estimated well).
``Regression'' here can mean anything from OLS to a random forest --- which
is exactly why cross-fitting (train/evaluate on separate splits) is needed
once flexible, ML-grade regressions are used for $\hat\mu$ and $\hat\pi$ and
a valid confidence interval is still wanted.

\subsection{A little machine learning}

$f$ is ``any learner'' --- linear model, random forest, neural net, whatever
was fit to predict cardiovascular risk. This project treats $f$ as a
\textbf{black box}, queried only by feeding it synthetic inputs and reading
off predictions (never inspecting its internals) --- that is what makes CDPs
``model-agnostic.''

Two ML facts turn into project rules, not just background color:
\begin{itemize}
  \item Flexible learners can nearly reproduce their own training outcomes
    (a forest is a classic case) --- this is \emph{why}
    \S\ref{sec:scope} insists explanatory curves are drawn on data $f$ never
    trained on.
  \item A learner that \textbf{splits on $A$} (e.g.\ a tree) can make the
    true dependence curve $\theta$ itself jumpy/discontinuous in $a$ --- this
    is \emph{why} \S\ref{sec:continuous} treats the smooth-curve estimator as
    needing an extra smoothness assumption a jumpy $\theta$ would violate.
\end{itemize}

\begin{openquestions}
  \item Work the toy backdoor-adjustment arithmetic above with actual
    made-up individual-level data (not just cell means) to see the
    confounded vs.\ adjusted contrast come out of a regression fit, not a
    hand calculation.
  \item Pin down precisely how positivity/overlap (potential-outcomes
    language) and the graph's positivity condition
    (\S\ref{sec:estimation}: ``the population positivity condition doesn't
    get harder'' when adjusting on more covariates) are the same statement
    in two vocabularies.
\end{openquestions}
